\section{The AI Native Coin}
\label{sec:aicoin}

\subsection{Token vs. coin}

LUX is the native coin of the Lux primary network and pays for base-layer
security and gas. AI is the native coin of Hanzo Network and rewards useful AI
work finalized through A-Chain PoAI consensus. They are independent assets;
neither is a wrapper of the other. Lux supplies the post-quantum root of
finality, while AI supplies the demand, staking, and reward economy for the
open compute network.

\subsection{Hard-cap allocation}

\begin{table}[h]
\centering
\small
\begin{tabular}{lr}
\toprule
\textbf{Parameter} & \textbf{Value} \\
\midrule
Maximum nominal supply & 2{,}000{,}000{,}000{,}000 AI \\
Permanent Hanzo DAO allocation & 1{,}000{,}000{,}000{,}000 AI \\
PoAI miner and validator pool & 1{,}000{,}000{,}000{,}000 AI \\
Destination-chain inflation & 0 AI \\
Decimals & 18 \\
\bottomrule
\end{tabular}
\caption{AI has one network-wide supply. Destination representations conserve
against canonical settlement and never add a local allocation.}
\label{tab:supply}
\end{table}

The DAO allocation is fixed at genesis and cannot be reclassified into mining
emission. It remains under the Hanzo DAO's constitutional treasury controls for
research, public infrastructure, ecosystem development, liquidity, and long-run
network stewardship. The proof-earned pool has no discretionary mint path: each
unit must fit the active era allowance and a finalized A-Chain mining receipt.

\subsection{Bitcoin-shaped proof-earned emission}

Let $M=10^{12}$ AI be the proof-earned pool, $h_0$ the activation height, and
$H$ the governance-fixed halving interval in finalized A-Chain blocks
(nominally four years). For era
$k=\lfloor(h-h_0)/H\rfloor$ and fractional progress $f\in[0,1)$, cumulative
mining allowance is

\begin{equation}
  A(h)=M(1-2^{-k})+fM2^{-(k+1)}.
\end{equation}

The first era may emit at most 500 billion AI, the second 250 billion, the
third 125 billion, and so on. Since
$\sum_{k\ge0}M2^{-(k+1)}=M$, miner and validator rewards converge to exactly
one trillion AI and total nominal issuance cannot exceed two trillion AI.

The canonical settlement state tracks \texttt{mintedMining} and clamps every
receipt to $A(h)-\texttt{mintedMining}$. A destination never evaluates the
schedule; it consumes the amount already authorized by A-Chain. The precise
within-era split between miners and validators, and the work-normalization
function across job classes, must be fixed and tested before activation.

\subsection{What earns AI}

AI is earned for proof-bearing useful work, not for owning hardware or
self-reporting time. A subsidy-eligible receipt must bind:

\begin{itemize}[leftmargin=1.2em]
  \item a pre-admitted requester or governed network-service job;
  \item a model, input commitment, deliverable, and consumer;
  \item a deterministic execution and proof policy;
  \item a committed output, transcript, and verifiable metering root;
  \item a post-commitment challenge or activated succinct proof;
  \item a unique A-Chain work nullifier and payout destination.
\end{itemize}

A raw GPU-hour, miner-created circular job, copied output, TEE quote, customer
signature, or destination-local proof is insufficient. This separation is what
lets the protocol reward actual work without turning fake invoices into a mint.

\subsection{Compute once, earn twice}

PoAI is an incremental revenue rail for existing AI clouds. A provider may keep
the ordinary customer fee for an inference or training workload and also earn a
declining AI reward by capturing proof evidence during that same execution. A
decentralized provider uses the identical interface to accept open jobs when
capacity would otherwise be idle.

The economic claim depends on verification asymmetry. For a supported matrix
operation the provider performs cubic work, while an ordinary CPU verifies a
post-commitment Freivalds opening in quadratic work plus commitment overhead.
Small jobs may be replayed exactly. Large jobs require graph coverage and data
availability under the optimistic policy, or an activated succinct proof.
Proof capture must remain materially cheaper than rerunning the workload.

\subsection{Fees, burns, and net deflation}

Demand fees and mining subsidy are separate. The requester pays for useful
service; the declining mining allowance bootstraps verifiable supply. A governed
fraction $\beta$ of settled protocol fees is burned:

\begin{equation}
  \Delta S = \text{PoAI issuance} - \beta\,\text{settled fees}.
\end{equation}

As geometric issuance tends toward zero, persistent fee burn can make net supply
deflationary. Burns do not replenish either the DAO allocation or the mining
pool.

\subsection{Staking and slashing}

Validators stake under the active Lux/A-Chain policy to order and finalize PoAI
state. Compute providers bond against objective fraud, transcript withholding,
metering fraud, replay, and policy violations. TEE collateral and measurement
rules apply only when a confidential job selects that evidence profile. Stake
requirements and reward weights are launch parameters, not immutable marketing
figures, and require measured economics and governance activation.
